% !TeX root = ../../Book2.tex
% 纲要标题：### 5.3 泛性质与普遍构造
% 纲要位置：工作记录/计划纲要.md。

\section{泛性质与普遍构造}
\label{b2:sec:0503}

同一个商模可以用不同生成元和关系表示，同一个直和也可以有不同的集合实现。若要比较这些实现，最有用的问题是：哪些映射能唯一地经过它？本节把这类泛性质写成统一的语言，并亲自构造模的拉回与推出。所有环仍保单位元，所有模均为左模；这些构造不要求环交换。

% 纲要内容（仅作写作计划，尚非教材正文）：
%
% * 初始对象、终对象、积与余积
% * 无限直积与指定嵌入不满足余积唯一性
% * 模的拉回与推出
% * 由 Hom 映射性质辨认对象
% * 正则模作为同构的测试对象
%

\subsection{初始对象、终对象、积与余积}
\label{b2:subsec:0503:01}

\begin{definition}\label{b2:def:initial-terminal-object}
设 \(\mathcal C\) 是范畴。若对象 \(I\) 到 \(\mathcal C\) 的每个对象 \(X\) 恰有一个态射，就称 \(I\) 为\term[chushi-duixiang]{初始对象}[initial object]。若每个 \(X\) 到对象 \(T\) 恰有一个态射，就称 \(T\) 为\term[zhong-duixiang]{终对象}[terminal object]。若一个对象兼为初始对象和终对象，就称它为\term[ling-duixiang]{零对象}[zero object]。
\end{definition}

集合范畴的初始对象为空集，终对象为任一单点集，故没有零对象；模范畴中的零模则为零对象。保单位元同态的环范畴中，\(\mathbb Z\) 是初始对象：到 \(R\) 的映射必为 \(n\mapsto n1_R\)，而此公式确实保持环运算与单位元。零环是终对象，却不是初始对象，因为从零环到非零环不存在保单位元同态。

\ProseParagraphBreak
两个初始对象 \(I,I'\) 之间的唯一态射 \(u:I\rightarrow I'\)、\(v:I'\rightarrow I\) 满足 \(vu=\operatorname{id}_I\)、\(uv=\operatorname{id}_{I'}\)，因为相应自态射集合各只有一个元素。因此初始对象之间有唯一同构；反转箭头即得终对象的相同结论。零对象的存在还为每对对象给出零态射，即经过零对象的唯一复合；在模范畴中它就是通常的零映射。

\begin{definition}\label{b2:def:product-coproduct}
设 \(\mathcal C\) 是范畴，\((X_i)_{i\in I}\) 为由集合 \(I\) 编号的 \(\mathcal C\) 中对象族。若对象 \(P\) 及态射 \(p_i:P\rightarrow X_i\) 满足：对任意对象 \(T\) 及态射族 \(a_i:T\rightarrow X_i\)，存在唯一态射 \(a:T\rightarrow P\) 使全部 \(p_i a=a_i\)，就称它们为该对象族的\term[ji]{积}[product]。

若对象 \(Q\) 及态射 \(j_i:X_i\rightarrow Q\) 满足：对任意对象 \(T\) 及态射族 \(b_i:X_i\rightarrow T\)，存在唯一态射 \(b:Q\rightarrow T\) 使全部 \(bj_i=b_i\)，就称它们为该对象族的\term[yuji]{余积}[coproduct]。
\end{definition}

这里的积是 \((P,(p_i)_{i\in I})\)，余积是 \((Q,(j_i)_{i\in I})\)：泛性质刻画的是对象及其指定的结构映射。例如，把 \(k^2\) 作为两个 \(k\) 的积时，两个坐标投影也是数据的一部分；只给出 \(k^2\) 的同构类型，还没有指定这项积构造。

\ProseParagraphBreak
积与余积的分解方向相反。下面两个三角形分别对每个 \(i\in I\) 成立；虚线是同时满足全族相容条件的唯一态射。
\[
\begin{tikzcd}[column sep=3em,row sep=2.5em]
T\arrow[r,dashed,"\exists!\,a"]\arrow[dr,"a_i"']&P\arrow[d,"p_i"]\\
&X_i
\end{tikzcd}
\qquad
\begin{tikzcd}[column sep=3em,row sep=2.5em]
X_i\arrow[r,"j_i"]\arrow[dr,"b_i"']&Q\arrow[d,dashed,"\exists!\,b"]\\
&T
\end{tikzcd}
\]
在模范畴中，积由逐坐标运算的 \(\prod_i M_i\) 及坐标投影给出。给定 \(a_i:T\rightarrow M_i\)，唯一可能的映射是 \(a(t)=\bigl(a_i(t)\bigr)_i\)，它逐坐标保持加法和左标量乘法，故确为模同态。余积由有限支集直和 \(\bigoplus_iM_i\) 及坐标嵌入给出。给定 \(b_i:M_i\rightarrow T\)，定义
\[
b\bigl((m_i)_i\bigr)=\sum_i b_i(m_i).
\]
有限支集保证右边有意义；分配律保证线性；每个直和元素都是有限个 \(j_i(m_i)\) 之和，保证唯一性。这正是\cref{b2:thm:module-product-universal}与\cref{b2:thm:module-sum-universal}在范畴语言中的表达。

\ProseParagraphBreak
有限个模的积与余积可以由同一个模承载，但其指定映射和泛性质方向不同。无限时不能混用：例如对域 \(k\)，直和 \(k^{(\mathbb N)}\) 中每个元素只有有限个非零坐标，而积 \(k^{\mathbb N}\) 还包含 \((1,1,\ldots)\)。恒等映射族 \(k\rightarrow k\) 合成的映射 \(k\rightarrow k^{\mathbb N}\)、\(a\mapsto(a,a,\ldots)\) 就不能落入直和。空积为终对象，空余积为初始对象，均可直接从定义读出。

\begin{proposition}\label{b2:prop:infinite-product-not-coproduct}
设 \(k\) 是域，\(P=k^{\mathbb N}\)，并令 \(j_n:k\to P\) 将一个标量送到仅第 \(n\) 个坐标取该值的向量。则 \(\bigl(P,(j_n)_{n\in\mathbb N}\bigr)\) 不是 \(k\) 向量空间范畴中这族 \(k\) 的余积。
\end{proposition}
\begin{proof}
令 \(S=k^{(\mathbb N)}\subseteq P\)，取 \(w=(1,1,\ldots)\notin S\)。在子空间 \(S+kw\) 中，表达 \(s+aw\) 唯一：若 \(aw\in S\) 且 \(a\ne0\)，则 \(aw\) 有无限多个非零坐标，矛盾。因此
\[
\ell_0(s+aw)=a
\]
定义线性泛函。由\cref{b2:lem:functional-extension}，它可以延拓为 \(P\) 上的线性泛函 \(\ell\)，满足 \(\ell(w)=1\) 且 \(\ell|_S=0\)。每个 \(j_n(k)\) 都在 \(S\) 中，故 \(\ell j_n=0\) 对全部 \(n\) 成立。然而零映射与 \(\ell\) 是两个不同的映射 \(P\to k\)，却在所有 \(j_n\) 上给出同一族零映射，违反余积的唯一性。
\end{proof}

\begin{proposition}\label{b2:prop:universal-object-uniqueness}
设 \((X_i)_{i\in I}\) 是范畴 \(\mathcal C\) 中由集合 \(I\) 编号的对象族。同一对象族的任意两个积之间，存在唯一保持全部投影的同构；两个余积之间，存在唯一保持全部指定嵌入的同构。
\end{proposition}
\begin{proof}
设 \((P,p_i)\)、\((P',p'_i)\) 都为积。用 \(P\) 及其映射族 \((p_i)\) 代入 \(P'\) 的泛性质，得唯一 \(u:P\rightarrow P'\) 满足 \(p'_iu=p_i\)；交换角色得 \(v:P'\rightarrow P\)。于是 \(p_ivu=p_i\)，而 \(\operatorname{id}_P\) 也满足这些等式。\(P\) 的唯一性迫使 \(vu=\operatorname{id}_P\)，同理 \(uv=\operatorname{id}_{P'}\)。\(u\) 的唯一性已经包含在构造中。

对余积 \((Q,j_i)\)、\((Q',j'_i)\)，两次泛性质分别给出 \(u j_i=j'_i\)、\(v j'_i=j_i\)。其复合在每个 \(j_i\) 上等于恒等映射，由余积的唯一性得 \(vu=\operatorname{id}_Q\)、\(uv=\operatorname{id}_{Q'}\)。
\end{proof}

“唯一”修饰保持指定结构映射的同构；抽去这些映射后，对象本身可能有许多自同构。例如 \(k^2\) 有许多可逆线性变换，只有恒等映射同时保持两个坐标投影。

\subsection{模的拉回与推出}
\label{b2:subsec:0503:02}

若两个模 \(M,N\) 都映到同一个模 \(P\)，就可以寻找像相等的成对元素：哪些 \((m,n)\) 满足 \(f(m)=g(n)\)？同样，从任意测试对象出发的两条映射也应在 \(P\) 中具有相同复合。拉回把这些相容的映射对合成一条映射。反过来，若两条映射 \(u:P\to M\)、\(v:P\to N\) 有公共来源，推出要给出一个共同目标，使来自同一 \(p\in P\) 的 \(u(p)\) 与 \(v(p)\) 在那里相等。这两种要求在一般范畴中分别表述如下。

\begin{definition}\label{b2:def:pullback-pushout}
设 \(\mathcal C\) 是范畴。给定态射 \(f:M\rightarrow P\)、\(g:N\rightarrow P\)。若对象 \(L\) 及态射 \(p:L\rightarrow M\)、\(q:L\rightarrow N\) 满足 \(fp=gq\)，并且对每对满足 \(fa=gb\) 的态射 \(a:T\rightarrow M\)、\(b:T\rightarrow N\)，存在唯一 \(h:T\rightarrow L\) 满足 \(ph=a\)、\(qh=b\)，就称它们为 \(f,g\) 的\term[lahui]{拉回}[pullback]。

给定态射 \(u:P\rightarrow M\)、\(v:P\rightarrow N\)。若对象 \(Q\) 及态射 \(i:M\rightarrow Q\)、\(j:N\rightarrow Q\) 满足 \(iu=jv\)，并且对每对满足 \(au=bv\) 的态射 \(a:M\rightarrow T\)、\(b:N\rightarrow T\)，存在唯一 \(h:Q\rightarrow T\) 满足 \(hi=a\)、\(hj=b\)，就称它们为 \(u,v\) 的\term[tuichu]{推出}[pushout]。
\end{definition}

四种构造都由测试对象的映射决定，区别在于箭头的方向以及拉回、推出所需的相容等式。下表中的 \(T\) 可为范畴中的任意对象。
\begin{table}[htbp]
\centering
\caption[积、余积、拉回与推出]{积、余积、拉回与推出的测试箭头}\label{b2:tab:universal-constructions}
\begin{tabularx}{\linewidth}{@{}p{0.13\linewidth}p{0.27\linewidth}X@{}}
\toprule
构造 & 给定图形 & 测试箭头与唯一分解 \\
\midrule
积 & \((X_i)_{i\in I}\) & 各 \(a_i:T\to X_i\) 唯一汇成 \(T\to\prod_iX_i\)。 \\
余积 & \((X_i)_{i\in I}\) & 各 \(b_i:X_i\to T\) 唯一汇成 \(\coprod_iX_i\to T\)。 \\
拉回 & \(M\xrightarrow{f}P\xleftarrow{g}N\) & \(a:T\to M\)、\(b:T\to N\) 满足 \(fa=gb\) 时，唯一汇成 \(T\to M\times_PN\)。 \\
推出 & \(M\xleftarrow{u}P\xrightarrow{v}N\) & \(a:M\to T\)、\(b:N\to T\) 满足 \(au=bv\) 时，唯一汇成 \(M\amalg_PN\to T\)。 \\
\bottomrule
\end{tabularx}
\end{table}
积与拉回接收从测试对象出发的箭头；余积与推出则把指向测试对象的箭头合成一条。拉回和推出还要求这两条箭头在公共对象处相容。

\ProseParagraphBreak
在模范畴中，拉回可以在 \(M\oplus N\) 内选择满足 \(f(m)=g(n)\) 的元素。推出则从同一份直和开始，把公共来源的两个像识别起来：直和内的 \((u(p),0)\) 与 \((0,v(p))\) 之差为 \((u(p),-v(p))\)，所以应将这些差值视为零。映射
\[
d:P\longrightarrow M\oplus N,\qquad
p\longmapsto\bigl(u(p),-v(p)\bigr)
\]
是模同态，它的像已是子模。因此拉回由满足约束的子模给出，推出由加入这些关系的商模给出。

\begin{proposition}\label{b2:prop:module-pullback-pushout}
设 \(R\) 是环，\(f:M\to P\)、\(g:N\to P\) 及 \(u:P\to M\)、\(v:P\to N\) 是左 \(R\) 模同态。在左 \(R\) 模范畴中，这两对同态的拉回和推出分别可取为
\begin{align}
M\times_PN&=\bigl\{\,(m,n)\in M\oplus N\mid f(m)=g(n)\,\bigr\},\label{b2:eq:module-pullback}\\
M\amalg_PN&=(M\oplus N)/\Bigl\{\,\bigl(u(p),-v(p)\bigr)\mid p\in P\,\Bigr\}.\label{b2:eq:module-pushout}
\end{align}
拉回的映射为坐标投影，推出的映射为 \(i(m)=\bigl[(m,0)\bigr]\)、\(j(n)=\bigl[(0,n)\bigr]\)。
\end{proposition}
\begin{proof}
若 \(f(m)=g(n)\)、\(f(m')=g(n')\)，则
\[
f(m+m')=g(n+n'),\qquad f(rm)=g(rn),
\]
所以式\eqref{b2:eq:module-pullback}定义子模，两个投影为模同态并满足相容条件。给定 \(a,b\) 且 \(fa=gb\)，令 \(h(t)=\bigl(a(t),b(t)\bigr)\)。相容条件保证其像在拉回中，坐标公式保证线性；两个投影的值又决定元素的两个坐标，所以 \(h\) 唯一。

式\eqref{b2:eq:module-pushout}的分母是同态 \(d:P\rightarrow M\oplus N\)、\(p\mapsto\bigl(u(p),-v(p)\bigr)\) 的像，因而是子模。商模存在，且
\[
iu(p)-jv(p)=\Bigl[\bigl(u(p),-v(p)\bigr)\Bigr]=0.
\]
若 \(au=bv\)，定义 \(s:M\oplus N\rightarrow T\)、\(s(m,n)=a(m)+b(n)\)。它线性，并在 \(d(P)\) 上为零，所以
\[
h\Bigl(\bigl[(m,n)\bigr]\Bigr)=a(m)+b(n)
\]
不依赖代表元，给出所需模同态。每个商元素都有形式 \(i(m)+j(n)\)，任何满足 \(hi=a\)、\(hj=b\) 的同态都必须有此公式，因而唯一。
\end{proof}

这些定义可用一对相反方向的交换图比较。外侧实线为任意满足定义中相容条件的 \(a,b\)，虚线为唯一的分解：
\[
\begin{tikzcd}[column sep=2.1em,row sep=2em]
T\arrow[rd,dashed,"\exists!\,h"]\arrow[rdd,bend right=14,"a"']\arrow[rrd,bend left=14,"b"]&&\\
&M\times_PN\arrow[r,"q"]\arrow[d,"p"']&N\arrow[d,"g"]\\
&M\arrow[r,"f"']&P
\end{tikzcd}
\qquad
\begin{tikzcd}[column sep=2.1em,row sep=2em]
P\arrow[r,"v"]\arrow[d,"u"']&N\arrow[d,"j"]\arrow[ddr,bend left=14,"b"]&\\
M\arrow[r,"i"']\arrow[drr,bend right=14,"a"']&M\amalg_PN\arrow[dr,dashed,"\exists!\,h"]&\\
&&T.
\end{tikzcd}
\]
交换性只是条件的一半，还必须验证任意相容映射的唯一分解。用两次泛性质构造比较态射，并用唯一性检验两个复合为恒等，就能像\cref{b2:prop:universal-object-uniqueness}一样证明拉回、推出在保持图中映射的同构意义下唯一。

\ProseParagraphBreak
例如在整数模中，\(f:\mathbb Z\rightarrow\mathbb Z\) 为乘二，\(g\) 为乘三。拉回的条件是 \(2m=3n\)，故 \(m=3t,n=2t\)，它与 \(\mathbb Z\) 同构，但两个投影分别为乘三、乘二。把两映射改读为从公共定义域出发的 \(u=2\)、\(v=3\)，推出为
\[
\mathbb Z^2/\mathbb Z(2,-3)\cong\mathbb Z,
\qquad \bigl[(m,n)\bigr]\longmapsto3m+2n.
\]
该映射满射，因为 \(3-2=1\)；其核由 \(3m=-2n\) 确定为 \(\mathbb Z(2,-3)\)，故\cref{b2:thm:module-first-isomorphism}给出同构。此时两个结构映射仍分别是乘三、乘二。同构类型相同并不消除它们所记录的原映射。

\ProseParagraphBreak
取零模到 \(P\) 的零映射，拉回可由 \(\ker f\) 及核嵌入 \(\iota\) 给出；取 \(P\) 到零模的零映射，推出可由 \(\operatorname{coker}u=M/u(P)\) 及商映射 \(q\) 给出。相应的交换图为
\[
\begin{tikzcd}[column sep=3em,row sep=2.2em]
\ker f\arrow[r,"\iota"]\arrow[d]&M\arrow[d,"f"]\\
0\arrow[r]&P
\end{tikzcd}
\qquad
\begin{tikzcd}[column sep=3em,row sep=2.2em]
P\arrow[r,"u"]\arrow[d]&M\arrow[d,"q"]\\
0\arrow[r]&\operatorname{coker}u.
\end{tikzcd}
\]
左图把 \(f(m)=0\) 作为子模条件，右图把 \(u(P)\) 中的元素作为要模掉的关系。核与余核因此是拉回、推出的特殊情形，也与\cref{b2:def:module-kernel-image-cokernel}中的具体定义一致。

\subsection{由 Hom 映射性质辨认对象}
\label{b2:subsec:0503:03}

前面通过泛性质证明了积、余积等构造在保持指定映射的同构意义下唯一。把它们的映射条件写成 Hom 集合之间的自然对应，还可以从映射性质辨认对象。以积 \((P,p_i)\) 为例，其泛性质给出双射
\[
\Phi_T:\operatorname{Hom}_{\mathcal C}(T,P)\longrightarrow
\prod_i\operatorname{Hom}_{\mathcal C}(T,X_i),\qquad
h\longmapsto(p_ih)_i.
\]
对 \(a:T'\rightarrow T\)，两侧都通过预复合 \(a\) 改变输入，得到交换方块
\[
\begin{tikzcd}[column sep=4em,row sep=2.5em]
\operatorname{Hom}_{\mathcal C}(T,P)
\arrow[r,"\Phi_T"]\arrow[d,"-\circ a"']&
\prod_i\operatorname{Hom}_{\mathcal C}(T,X_i)\arrow[d,"-\circ a"]\\
\operatorname{Hom}_{\mathcal C}(T',P)\arrow[r,"\Phi_{T'}"']&
\prod_i\operatorname{Hom}_{\mathcal C}(T',X_i).
\end{tikzcd}
\]
右侧的预复合逐坐标进行，两条路径都把 \(h\) 送到 \((p_iha)_i\)，所以这些双射关于 \(T\) 自然。拉回则把右端换成满足 \(fa=gb\) 的映射对；推出的对应从 \(\operatorname{Hom}_{\mathcal C}(Q,T)\) 出发，关于 \(T\) 协变。它们不只数出了映射有多少，而是给出了与复合相容的具体对应。

\begin{proposition}[由映射检测同构]\label{b2:prop:hom-detects-isomorphism}
设 \(\mathcal C\) 是范畴，\(u:P\rightarrow Q\) 为 \(\mathcal C\) 中的态射。若对每个对象 \(T\)，后复合映射
\[
u_*:\operatorname{Hom}_{\mathcal C}(T,P)\longrightarrow
\operatorname{Hom}_{\mathcal C}(T,Q),\qquad h\longmapsto uh
\]
都是双射，则 \(u\) 为同构。更一般地，任何自然变换
\(\alpha:\operatorname{Hom}_{\mathcal C}(-,P)\Rightarrow
\operatorname{Hom}_{\mathcal C}(-,Q)\)
都唯一地由某个态射 \(u:P\rightarrow Q\) 的后复合给出；这个自然变换为自然同构，当且仅当 \(u\) 为同构。
\end{proposition}
\begin{proof}
先取 \(T=Q\)。满射性给出 \(v:Q\rightarrow P\)，使 \(uv=\operatorname{id}_Q\)。再取 \(T=P\)，有 \(u(vu)=u=u\operatorname{id}_P\)；这次的单射性给出 \(vu=\operatorname{id}_P\)，所以 \(u\) 可逆。

对自然变换 \(\alpha\)，若它由某个 \(u:P\to Q\) 后复合给出，在 \(T=P\)、\(h=\operatorname{id}_P\) 处就必有 \(\alpha_P(\operatorname{id}_P)=u\operatorname{id}_P=u\)。因此应当用恒等态射的像恢复候选态射，令 \(u=\alpha_P(\operatorname{id}_P)\)。给定 \(h:T\rightarrow P\)，用 \(h\) 的预复合及 \(\alpha\) 的自然性，得到
\[
\alpha_T(h)=\alpha_T(\operatorname{id}_P h)
=\alpha_P(\operatorname{id}_P)h=uh.
\]
所以 \(\alpha\) 的全部分量都由 \(u\) 后复合给出，在 \(T=P\) 的恒等态射处取值又说明 \(u\) 唯一。若 \(\alpha\) 是自然同构，其分量都为双射，第一部分说明 \(u\) 是同构。反过来，若 \(u\) 是同构，后复合 \(u^{-1}\) 就给出每个分量的逆，所以 \(\alpha\) 是自然同构。
\end{proof}

\begin{corollary}[正则模检测同构]\label{b2:cor:regular-module-detects}
设 \(R\) 是环，\(f:M\to N\) 是左 \(R\) 模同态。后复合映射
\[
\operatorname{Hom}_R(R,f):\operatorname{Hom}_R(R,M)
\longrightarrow\operatorname{Hom}_R(R,N)
\]
为双射，当且仅当 \(f\) 是模同构。
\end{corollary}
\begin{proof}
由\cref{b2:prop:regular-hom-endomorphism}，求值 \(\varphi\mapsto\varphi(1_R)\) 分别将两个 Hom 群识别为 \(M,N\)。而 \((f\varphi)(1_R)=f(\varphi(1_R))\)，所以后复合映射在这些求值对应下就是 \(f\) 本身。若它双射，\(f\) 是双射模同态，因而是\cref{b2:def:module-homomorphism}所定义的模同构；逆映射保持模运算。反过来，后复合 \(f^{-1}\) 给出 \(\operatorname{Hom}_R(R,f)\) 的逆。
\end{proof}

一般固定测试对象并没有正则模的这种能力。例如，包含同态 \(\mathbb Z\to\mathbb Q\) 不是同构，但它在 \(\operatorname{Hom}_{\mathbb Z}(\mathbb Z/2\mathbb Z,-)\) 下成为两个零群之间的双射：从 \(\mathbb Z/2\mathbb Z\) 出发的同态，其生成元像必须被 \(2\) 零化，而 \(\mathbb Z,\mathbb Q\) 中满足这一条件的都只有零。因此“对每个 \(T\)”不能无条件换成任意一个固定 \(T\)。

\ProseParagraphBreak
第四卷的 Yoneda 引理将把上述第二个 Hom 函子换成一般集合值反变函子，以同样的恒等态射求值思想描述自然变换。这里要求的是与所有复合相容的具体对应；即使各个 Hom 集合大小恰好相同，也没有自动得到双射的自然相容性。
